% !TEX program = xelatex
\documentclass[11pt, a4paper]{article}

\usepackage[margin=2.5cm]{geometry}
\usepackage{fontspec}
\usepackage{booktabs}
\usepackage{longtable}
\usepackage{array}
\usepackage{hyperref}
\usepackage{parskip}
\usepackage{enumitem}

\setmainfont{Times New Roman}
\setmonofont{Courier New}

\hypersetup{colorlinks=false, pdftitle={Output Dataset Codebook}}

\newcolumntype{P}[1]{>{\raggedright\arraybackslash}p{#1}}
\newcommand{\ra}[1]{\renewcommand{\arraystretch}{#1}}
\newcommand{\thead}[1]{\textbf{#1}}

% Quoted prompt text block: indented, slightly smaller, used to reproduce
% verbatim system/user prompt wording without switching to a monospace code font.
\newenvironment{promptquote}{%
  \begin{quote}\small
}{%
  \end{quote}
}

% ═══════════════════════════════════════════════════════════════════════════════
\begin{document}
% ═══════════════════════════════════════════════════════════════════════════════

\begin{center}
  {\LARGE\bfseries The Growth of Executive Constraints}\\[0.4cm]
  {\large LLM Pipeline Codebook}\\[0.2cm]
  {Peng-Ting Kuo} \\[0.2cm]
  {\small \today}
\end{center}

\bigskip

This codebook documents the dataset produced by the Gemini Batch API pipeline.
It covers every column in the canonical, full-dataset output files. The dataset can be found in the \href{https://drive.google.com/drive/folders/1Y1oUJ41vCwXSeHYVmyZUn2sG8uHMwJx6?usp=sharing}{Google Drive folder}.

\medskip
\textbf{Canonical run settings.} The documented dataset is built WITHOUT
self-consistency verification and WITHOUT reasoning. Every
indicator therefore contributes exactly two plain columns —
\texttt{\{ind\}\_\allowbreak prediction} and \texttt{\{ind\}\_\allowbreak confidence} —
except \texttt{constitution} (three extra document-level columns) and
\texttt{elections} (a downstream pass-through rule).

% ═══════════════════════════════════════════════════════════════════════════════
\section{Input Columns}
% ═══════════════════════════════════════════════════════════════════════════════

The pipeline input is a leader-level table; the following columns from the original dataset are
used.

\ra{1.3}
\begin{longtable}{P{3.8cm} P{1.5cm} P{9cm}}
  \toprule
  \thead{Column} & \thead{Type} & \thead{Meaning} \\
  \midrule
  \endhead
  \bottomrule
  \endfoot
  \texttt{id}                 & string/int & Unique row identifier \emph{[originated from
    \texttt{id} column]} \\
  \texttt{slug\_id}            & string     & Informative composite key, \texttt{polity\_\allowbreak
    leader\_\allowbreak years}; not guaranteed unique in the raw source \emph{[originated from
    \texttt{slug\_id} column]} \\
  \texttt{polity\_name}        & string     & Name of the historical polity \emph{[originated
    from \texttt{territory\allowbreak name\allowbreak historical} column]} \\
  \texttt{leader\_name}        & string     & Name of the leader \emph{[originated from
    \texttt{name\_\allowbreak clean} column]} \\
  \texttt{leader\_first\_year} & Int64 (nullable) & First year of tenure, negative = BCE
    \emph{[originated from \texttt{entry\allowbreak date\allowbreak year} column]} \\
  \texttt{leader\_last\_year}  & Int64 (nullable) & Last year of tenure, negative = BCE
    \emph{[originated from \texttt{exit\allowbreak date\allowbreak year} column]} \\
\end{longtable}

Missing \texttt{leader\_\allowbreak first\_\allowbreak year} / \texttt{leader\_\allowbreak
last\_\allowbreak year} values are rendered as the literal string \\ \texttt{"unknown"} inside
prompts.

% ═══════════════════════════════════════════════════════════════════════════════
\section{Indicator Definitions and Output Columns}
% ═══════════════════════════════════════════════════════════════════════════════

Predictions are produced by three separate tasks that are later merged on \texttt{id}:
\texttt{constitution} (its own dedicated prompt), \texttt{indicators} (all 17 non-constitution,
non-elections indicators, combined into ONE prompt per leader-tenure), and \texttt{elections}
(downstream, depends on \texttt{assembly\_\allowbreak prediction}). Every column pair
below is \texttt{\{ind\}\_\allowbreak prediction} (the coded label) and
\texttt{\{ind\}\_\allowbreak confidence} (LLM self-reported confidence, integer 1--100).


% ─────────────────────────────────────────────────────────────────────────────
\subsection{Constitution (\texttt{constitution})}
% ─────────────────────────────────────────────────────────────────────────────

\textbf{Prompt source:} \texttt{prompts/\allowbreak constitution.py}. \\
\textbf{Clarficiation:} Constitution is never part of the combined \texttt{indicators} prompt — 
it always runs as its own task/LLM call, using CoT and few-shot prompt.

\begin{promptquote}
You are a professional political scientist, historian, and constitutional expert specializing
in constitutional history across different countries.

Your task is to classify the HIGHEST TYPE of written legal document that existed in a given
polity during a specific leader's tenure.

\textbf{Type 0 --- No Written Legal Document.} No written code of law or constitution existed
during this period. Only oral traditions, customary practices, or uncodified norms are present.

\textbf{Type 1 --- Code of Law (Subjects/Citizens Only).} A written document containing binding
legal provisions that govern the conduct of subjects or citizens, but does NOT constrain rulers
or specify governance rules. Both conditions must be met: (1) Written Document --- an
identifiable written document with legal force; (2) Code of Law --- binding legal rules
governing the conduct of subjects/citizens. Examples: Code of Ur-Nammu in Sumer
(c.\ 2100--2050 BCE), Code of Hammurabi in Babylon (c.\ 1750 BCE), Law of Moses (Torah) in the
Kingdom of the Jews, Hittite Laws (Code of the Nesilim) in the Hittite Kingdom
(c.\ 1650--1500 BCE). Does NOT qualify: royal decrees or edicts (not a systematic code of law);
purely oral traditions or uncodified customs; documents that also specify governance rules or
limit rulers (those qualify for Type 2 instead).

\textbf{Type 2 --- Constitution (Code of Law + Governance Rules).} All requirements for Type 1
plus the document also governs the polity itself. All four elements must be present:
(1) Written Document --- an identifiable written document;
(2) Code of Law --- binding legal rules with legal force;
(3) Rules About Governance --- specifies how leaders are selected, how laws are enacted, or the
structure of governmental institutions;
(4) Limitations on Authority --- constrains the ruler's or government's power, through rights of
subjects, checks and balances, or procedures limiting arbitrary power.
Examples: Draconian constitution in Athens (late 7th c.\ BCE), Solonian constitution in Athens
(early 6th c.\ BCE), Law of Manu (Manusmriti) in India (c.\ 200 BCE) [borderline], Magna Carta in
England (1215), Golden Bull in Hungary (1222), Salic Law in the Frankish kingdoms (c.\ 500 AD).

\textbf{Classification logic:} (1) Is there a written document with binding legal force applying
to subjects? If NO $\rightarrow$ Type 0. (2) Does it ONLY govern subjects without constraining
rulers or specifying governance? $\rightarrow$ Type 1. (3) Does it ALSO specify governance rules
AND limit ruler authority? $\rightarrow$ Type 2.

\textbf{If multiple documents existed during the tenure:} classify each document separately;
report the highest type as \texttt{constitution}; list ALL documents in \texttt{document\_name}
(semicolon-separated); list their corresponding types in \texttt{document\_types}
(semicolon-separated, same order); list their adoption years in \texttt{constitution\_year}
(semicolon-separated, same order).

\textbf{Analysis process:}
\begin{enumerate}[leftmargin=1.4em]
  \item \textbf{Identify the Historical Context} --- What type of polity was this during this
    leader's tenure? Was it independent or subject to external governance authority?
  \item \textbf{Determine Governance Authority} --- Who had the authority to establish legal
    rules during this tenure? Was governance authority internal or externally imposed?
  \item \textbf{Identify Written Legal Documents} --- What specific written document(s) existed
    during this leader's tenure? Name each document with its official title and adoption year.
    List ALL documents, not only the highest-type one.
  \item \textbf{Classify Each Document} --- For each document: Does it have legal force
    (required for Type 1 and 2)? Does it apply to subjects/citizens (required for Type 1 and 2)?
    Does it specify governance rules (required for Type 2 only)? Does it limit the ruler's
    authority or protect subjects' rights (required for Type 2 only)?
  \item \textbf{Consider Temporal Factors} --- Did documents exist for the entire tenure, or only
    part of it? Were there amendments, new documents, or abrogations during this tenure?
  \item \textbf{Assess Confidence} --- How well-documented is the evidence? Is there scholarly
    consensus or debate? When uncertain, lean toward a lower type (conservative coding).
\end{enumerate}

\textbf{Confidence score guidelines:} 81--100 specific documents named with exact dates, all
elements clearly documented, scholarly consensus; 61--80 strong documentary evidence, minor
scholarly debate possible; 41--60 documents identified but some elements ambiguous, reasonable
alternative interpretations; 21--40 limited evidence, educated inference from general
regional/period knowledge; 1--20 no reliable specific information, essentially speculative.

\end{promptquote}

\medskip
\textit{User prompt template:}
\begin{quote}\small
Classify the written legal documents for the following leader's tenure: \\ 
\textbf{Polity:} \{polity\} \\
\textbf{Leader:} \{name\} \\
\textbf{Tenure Period:} \{start\_year\}--\{end\_year\}. \\
Work through the six analysis steps and return the JSON object specified above.  \\
When uncertain, default to a lower type (conservative coding).
\end{quote}

\textbf{Related columns:}
\ra{1.3}
\begin{longtable}{P{4.6cm} P{2cm} P{8.4cm}}
  \toprule
  \thead{Column} & \thead{Type} & \thead{Meaning} \\
  \midrule
  \endhead
  \bottomrule
  \endfoot
  \texttt{constitution\_\allowbreak prediction} & int (0/1/2) &
    The highest type of written legal document that existed in the polity during this leader's tenure. Codes:\newline
    0 = no written law\newline
    1 = code of law (governs subjects only)\newline
    2 = full constitution (written law + governance rules + limits on authority) \\
  \texttt{constitution\_\allowbreak confidence} & int (1--100) & LLM self-reported confidence score \\
  \texttt{constitution\_\allowbreak document\_\allowbreak name} & string      & Semicolon-separated name(s) of the
    referenced document(s), or \texttt{"N/A"} \\
  \texttt{constitution\_\allowbreak year} & string      & Semicolon-separated adoption year(s),
    same order as \texttt{document\_\allowbreak name} \\
  \texttt{constitution\_\allowbreak document\_\allowbreak types} & string      & Semicolon-separated type code(s)
    (1 or 2), same order as \texttt{document\_\allowbreak name} \\
\end{longtable}

% ─────────────────────────────────────────────────────────────────────────────
\subsection{Indicators (combined \texttt{indicators} task)}
% ─────────────────────────────────────────────────────────────────────────────

\textbf{Prompt source:} \texttt{prompts/single\_builder.py}. Sovereign, federalism, the nine
\texttt{checks\_*} sub-indicators, collegiality, petition, assembly, entry, exit, and symbolism
are NOT queried one prompt per indicator --- they are combined into a single LLM call per
leader-tenure.

The shared default system prompt opens with:

\begin{promptquote}
You are a political scientist coding executive constraints for historical leaders. Core rule:
focus on THIS specific leader's tenure. When evidence is uncertain, apply the
indicator-appropriate default: institutional-presence indicators (federalism, checks,
collegiality, petition, legislative election) default to 0 / None / No --- if such an
institution existed, the historical record would usually mention it, so silence indicates
absence. Sovereignty defaults to 1 / Sovereign --- silence indicates the polity governed its own
domestic affairs (be more cautious for premodern and non-Western polities, where semi-sovereign
status may go unrecorded). Assembly is ordinal (0$<$1$<$2$<$3): when choosing among
present-but-ambiguous levels, prefer the lower level. Nominal indicators (entry, exit) and the
non-monotonic symbolism scale do NOT default to a lower label --- lower the confidence score
instead of forcing a code.
\end{promptquote}

\ldots followed by the ``Indicator Definitions'' section reproduced indicator-by-indicator
below, and closing with a single JSON-output instruction requesting, for every selected
indicator, its coded label and \texttt{\{ind\}\_\allowbreak confidence} (1--100 integer).

\bigskip
\textbf{Sovereign (\texttt{sovereign})}

\begin{promptquote}
Sovereignty refers to a polity's ability to conduct its own domestic affairs without foreign
interference. Code the polity's de facto control over domestic affairs during this leader's
tenure; the conduct of foreign affairs is not relevant here. IMPORTANT: when evidence is
uncertain, default to 1 (Sovereign) --- silence usually means no foreign overlord --- but be
cautious for premodern/non-Western polities where tributary or vassal status may go unrecorded.

Coding: 0 = Semi-sovereign (colony, protectorate, distant or overseas territory not fully
incorporated into the metropole); 1 = Sovereign (city-states, nation-states, empires, republics,
monarchies, tributary states so long as they held primary responsibility for domestic affairs,
states within the Peloponnesian League, the Hanseatic League, the Holy Roman Empire, and the
European Union, any state recognized by European powers or international law in the modern era,
as well as a few that enjoy de facto sovereignty such as Somaliland and Taiwan).
\end{promptquote}

\texttt{sovereign\_prediction} (int, 0--1):
\begin{itemize}[nosep, leftmargin=1.7em]
  \item 0 = semi-sovereign
  \item 1 = sovereign domestic governance
\end{itemize}
\texttt{sovereign\_confidence} (int, 1--100): LLM self-reported confidence score.

\bigskip
\textbf{Federalism (\texttt{federalism})}

\begin{promptquote}
Federalism refers to a division of sovereignty between central and local units, reserving some
important powers to the latter and providing a relatively decentralized mode of governance. Such
powers are typically enshrined in a constitution, and localities are often represented in a
chamber at the polity level.

Coding: 0 = Non-federal; 1 = Federal --- includes confederations, leagues, and composite
monarchies. Historical examples: Achaean League, Aetolian League, Lycian League, Boeotian
League, Old Swiss Confederacy, Dutch Republic, Holy Roman Empire, Iroquois Confederacy,
Hanseatic League, Polish-Lithuanian Commonwealth, Tokugawa Japan. Contemporary examples: Canada,
Germany, India, United States, and the European Union.
\end{promptquote}

\texttt{federalism\_prediction} (int, 0--1):
\begin{itemize}[nosep, leftmargin=1.7em]
  \item 0 = unitary state
  \item 1 = federal --- sovereignty divided between central and local units
\end{itemize}
\texttt{federalism\_confidence} (int, 1--100): LLM self-reported confidence score.

\bigskip
\textbf{Collegiality (\texttt{collegiality})}

\begin{promptquote}
Collegiality refers to power at the apex of a polity being exercised in a shared manner:
decisionmaking power is divided among co-equal actors. There may be a titular head (e.g., a
director or chair), but the other members are regarded as partners and collaborators;
decisionmaking is lateral rather than vertical. Code de facto practice --- if a body is formally
collegial but actually dominated by a single actor, it is NOT collegial.

Coding: 0 = Non-collegial (most presidencies, monarchies, and dictatorships); 1 = Collegial (most
cabinets, some military juntas, many regencies, Roman consuls, Switzerland's modern presidency).
\end{promptquote}

\texttt{collegiality\_prediction} (int, 0--1):
\begin{itemize}[nosep, leftmargin=1.7em]
  \item 0 = single actor dominates decision-making
  \item 1 = decisions genuinely shared by a constituted body
\end{itemize}
\texttt{collegiality\_confidence} (int, 1--100): LLM self-reported confidence score.

\bigskip
\textbf{Petition (\texttt{petition})}

\begin{promptquote}
A petition is a regularized, institutionalized process by which a subject or citizen may lodge a
complaint or request for redress with a high official (e.g., a head of state, legislature,
court, or ombudsman). It may take the form of a face-to-face meeting (the ``bell of justice''
tradition), a letter, or an electronic communication, with one or many signatories.

Coding: 0 = No --- use of petition is extremely rare and probably ineffective, or there is no
record of its existence; 1 = Yes --- petitions are a fairly regular feature of political life.
\end{promptquote}

\texttt{petition\_prediction} (int, 0--1):
\begin{itemize}[nosep, leftmargin=1.7em]
  \item 0 = petitioning is rare, ineffective, or unrecorded
  \item 1 = petitions are a regular, institutionalized feature of political life
\end{itemize}
\texttt{petition\_confidence} (int, 1--100): LLM self-reported confidence score.

\bigskip
\textbf{Assembly (\texttt{assembly})}

\begin{promptquote}
An assembly is a body designed to govern directly, to select leaders, or to assist in governing.
It may be advisory (a council of selected elites), representative (a legislature), or inclusive
of all citizens (a popular assembly).

Coding: 0 = None. 1 = Council --- a small advisory council appointed by the ruler, which may not
enjoy much autonomy but is nonetheless institutionalized (has a recognized name, meets regularly,
and has a fairly stable membership); examples: noble or aristocratic councils, privy councils,
dynastic councils, Ottoman divan. 2 = Legislature --- a large representative body that plays some
role in policymaking or leadership selection, de jure or de facto; examples: estates assemblies
in premodern Europe, the Hwabaek Council in Korea during the Silla Dynasty, legislatures in
modern governments. 3 = Popular assembly --- an assembly that includes most citizens of the
polity or a representative sample chosen by lot; examples: Ecclesia in ancient Athens,
Landsgemeinden in modern Swiss cantons.
\end{promptquote}

\texttt{assembly\_prediction} (int, 0--3):
\begin{itemize}[nosep, leftmargin=1.7em]
  \item 0 = none
  \item 1 = council --- small advisory council (institutionalized)
  \item 2 = legislature --- policymaking or leadership-selection role
  \item 3 = popular assembly --- most citizens, or a representative sample chosen by lot
\end{itemize}
\texttt{assembly\_confidence} (int, 1--100): LLM self-reported confidence score.

\bigskip
\textbf{Entry (\texttt{entry})}

\begin{promptquote}
IMPORTANT: Code the manner of entry at the START of this tenure, not later events. If it cannot
be determined, use 99 (do NOT guess to avoid it). The manner in which the executive entered
office. Code the manner of entry at the start of this leader's tenure (the accession in
\{start\_year\}).

Coding: 0 = through the threat or application of force (coup, rebellion); 1 = appointed by a
foreign power; 2 = appointed by the ruling party (one-party system); 3 = appointed by a royal
council (royal family or conclave of aristocrats); 4 = through hereditary succession; 5 =
appointed by the military; 6 = appointed by the legislature; 7 = appointed by the head of state;
8 = appointed by the head of government; 9 = directly through a popular election (regardless of
the extension of the suffrage); 10 = other (including clerical bodies such as the College of
Cardinals); 99 = unknown (circumstances of entry are unknown).
\end{promptquote}

\texttt{entry\_prediction} (int, 0--10 or 99) --- how the executive came to power:
\begin{itemize}[nosep, leftmargin=1.7em]
  \item 0 = force (coup, rebellion, conquest)
  \item 1 = foreign power appointment
  \item 2 = ruling party appointment (one-party system)
  \item 3 = royal council appointment
  \item 4 = hereditary succession
  \item 5 = military appointment
  \item 6 = legislature appointment
  \item 7 = head of state appointment
  \item 8 = head of government appointment
  \item 9 = direct popular election
  \item 10 = other (e.g., clerical bodies such as the College of Cardinals)
  \item 99 = unknown
\end{itemize}
\texttt{entry\_confidence} (int, 1--100): LLM self-reported confidence score.

\bigskip
\textbf{Exit (\texttt{exit})}

\begin{promptquote}
IMPORTANT: Code the manner of exit at the END of this tenure, not later events. If it cannot be
determined, use 99 (do NOT guess to avoid it). The circumstances of the executive's departure
from office. Code the manner of exit at the end of this leader's tenure (the departure in
\{end\_year\}).

Coding: 0 = abdicated or retired voluntarily but NOT due to ill health; 1 = other regular exit
(e.g., term limits or defeat in election); 2 = transition to another office type/typology (by
regular procedures); 3 = died (of disease or accident) on campaign in civil war; 4 = died (of
disease or accident) on campaign in foreign war; 5 = died of other natural causes; 6 = retired
due to ill health; 7 = suicide; 8 = deposed by domestic actors; 9 = assassinated or forced
suicide; 10 = died in battle in civil war; 11 = died in battle in foreign war; 12 = transition to
another office by irregular procedures; 13 = through deposition by a foreign state; 14 = still in
office; 99 = unknown.
\end{promptquote}

\texttt{exit\_prediction} (int, 0--14 or 99) --- circumstances of the executive's departure:
\begin{itemize}[nosep, leftmargin=1.7em]
  \item 0 = voluntary retirement or abdication (not due to ill health)
  \item 1 = other regular exit (term limits, electoral defeat)
  \item 2 = regular transition to another office
  \item 3 = died on campaign in civil war (disease or accident)
  \item 4 = died on campaign in foreign war (disease or accident)
  \item 5 = died of natural causes
  \item 6 = retired due to ill health
  \item 7 = suicide
  \item 8 = deposed by domestic actors
  \item 9 = assassinated or forced suicide
  \item 10 = killed in battle in civil war
  \item 11 = killed in battle in foreign war
  \item 12 = irregular transition to another office
  \item 13 = deposed by a foreign state
  \item 14 = still in office at end of observation period
  \item 99 = unknown
\end{itemize}
\texttt{exit\_confidence} (int, 1--100): LLM self-reported confidence score.

\bigskip
\textbf{Symbolism (\texttt{symbolism})}

\begin{promptquote}
IMPORTANT: this scale is NON-MONOTONIC with respect to leader power --- a higher code does NOT
mean more actual power. Codes 0--2 rise with the grandeur of the office, but code 3 marks largely
ceremonial figureheads with little real power. Code the trappings and self-presentation of the
office, NOT the leader's actual power. (Purely ceremonial heads who are not paramount leaders are
excluded from the sample.)

Symbolic power is reflected in the trappings of office and the presentation of self --- e.g., a
grandly appointed palace; regalia such as scepters, seals, thrones, and special garments; an
aristocratic court and retinue; performance of spiritual rituals central to the polity; powers
normally reserved for deities; special forms of address marking the ruler's apartness; and
protections of the ruler's status such as l\`ese-majest\'e.

Coding: 0 = Plain --- the trappings of the office are plain and simple, little distinguishes the
personage of the ruler from others in the realm (example: British prime minister, 10 Downing
Street). 1 = Decorated --- the trappings of the office are impressive but connected to the office
rather than the officeholder, who is understood as mortal (example: American president). 2 =
Deified --- the trappings of the office are impressive and the holder is regarded as having
divine or quasi-divine status, separate and apart from mere mortals (example: many kings in the
premodern era). 3 = Ceremonial --- the trappings of office are so extensive, and so demanding,
that they serve as a constraint on the exercise of power, separating the leader from the springs
of policymaking; the executive's role is as much ceremonial as executive, approval of initiatives
is formal (example: Japanese emperor during most periods).
\end{promptquote}

\texttt{symbolism\_prediction} (int, 0--3; non-monotonic):
\begin{itemize}[nosep, leftmargin=1.7em]
  \item 0 = plain --- little distinguishes the ruler from others in the realm
  \item 1 = decorated --- impressive, but the officeholder is understood as mortal
  \item 2 = deified --- the holder is regarded as divine or quasi-divine
  \item 3 = ceremonial --- trappings so extensive they constrain the exercise of power
\end{itemize}
\texttt{symbolism\_confidence} (int, 1--100): LLM self-reported confidence score.

% ─────────────────────────────────────────────────────────────────────────────
\subsection{Checks --- Nine Independent Binary Sub-Indicators}
% ─────────────────────────────────────────────────────────────────────────────

The former 0--9 multi-select \texttt{checks} indicator was replaced with nine independent
binary questions, one per actor, so that self-consistency agreement is well-defined per actor.
``None'' is not itself a question --- it is the case where all nine columns equal 0. Each of the
nine shares the same framing and coding scale (independence + de facto capacity to resist the
executive; media, civil society, and ordinary citizens are NOT counted), reproduced below
indicator-by-indicator like the other core indicators.

\bigskip
\textbf{Checks: Local / Constituent Units (\texttt{checks\_local})}

\begin{promptquote}
Do local or constituent territorial/kin-based units provide an effective check on the executive
during THIS leader's tenure? An effective check means the body is independent of the executive
and has the de facto capacity to resist or constrain executive action --- not merely formal/de
jure standing. A body that exists on paper but is controlled by the executive does NOT count.
Media, civil society organizations, and ordinary citizens are NOT counted.

Examples: regional and local governments, tribes, clans, ethnic governance units.

Coding: 0 = No --- such units do not exist during the tenure, or exist but cannot effectively
resist the executive. 1 = Yes --- such units are independent and can, at least on occasion,
effectively resist the executive.
\end{promptquote}

\texttt{checks\_local\_prediction} (int, 0--1):
\begin{itemize}[nosep, leftmargin=1.7em]
  \item 0 = does not effectively check the executive
  \item 1 = effectively checks the executive
\end{itemize}
\texttt{checks\_local\_confidence} (int, 1--100): LLM self-reported confidence score.

\bigskip
\textbf{Checks: Military (\texttt{checks\_military})}

\begin{promptquote}
Do the armed forces or a warrior estate provide an effective check on the executive during THIS
leader's tenure? An effective check means the body is independent of the executive and has the
de facto capacity to resist or constrain executive action --- not merely formal/de jure standing.
A body that exists on paper but is controlled by the executive does NOT count. Media, civil
society organizations, and ordinary citizens are NOT counted.

Examples: officers, specific branches of the military, a warrior caste (e.g., Samurai).

Coding: 0 = No --- the military does not exist as an independent force, or cannot effectively
resist the executive. 1 = Yes --- the military is independent and can, at least on occasion,
effectively resist the executive.
\end{promptquote}

\texttt{checks\_military\_prediction} (int, 0--1):
\begin{itemize}[nosep, leftmargin=1.7em]
  \item 0 = does not effectively check the executive
  \item 1 = effectively checks the executive
\end{itemize}
\texttt{checks\_military\_confidence} (int, 1--100): LLM self-reported confidence score.

\bigskip
\textbf{Checks: Clergy (\texttt{checks\_clergy})}

\begin{promptquote}
Do religious authorities provide an effective check on the executive during THIS leader's
tenure? An effective check means the body is independent of the executive and has the de facto
capacity to resist or constrain executive action --- not merely formal/de jure standing. A body
that exists on paper but is controlled by the executive does NOT count. Media, civil society
organizations, and ordinary citizens are NOT counted.

Clergy power derives from their role as arbiters of a widely espoused religion or belief system.
Examples: established church, priesthood, religious caste (of any religion or denomination).

Coding: 0 = No --- no independent religious authority exists, or it cannot effectively resist
the executive. 1 = Yes --- religious authorities are independent and can, at least on occasion,
effectively resist the executive.
\end{promptquote}

\texttt{checks\_clergy\_prediction} (int, 0--1):
\begin{itemize}[nosep, leftmargin=1.7em]
  \item 0 = does not effectively check the executive
  \item 1 = effectively checks the executive
\end{itemize}
\texttt{checks\_clergy\_confidence} (int, 1--100): LLM self-reported confidence score.

\bigskip
\textbf{Checks: Aristocracy (\texttt{checks\_aristocracy})}

\begin{promptquote}
Does a hereditary or titled upper stratum provide an effective check on the executive during
THIS leader's tenure? An effective check means the body is independent of the executive and has
the de facto capacity to resist or constrain executive action --- not merely formal/de jure
standing. A body that exists on paper but is controlled by the executive does NOT count. Media,
civil society organizations, and ordinary citizens are NOT counted.

Examples: landed class, nobility, hereditary elite, titled class, patriciate, upper caste.

Coding: 0 = No --- no independent aristocracy exists, or it cannot effectively resist the
executive. 1 = Yes --- the aristocracy is independent and can, at least on occasion, effectively
resist the executive.
\end{promptquote}

\texttt{checks\_aristocracy\_prediction} (int, 0--1):
\begin{itemize}[nosep, leftmargin=1.7em]
  \item 0 = does not effectively check the executive
  \item 1 = effectively checks the executive
\end{itemize}
\texttt{checks\_aristocracy\_confidence} (int, 1--100): LLM self-reported confidence score.

\bigskip
\textbf{Checks: Bourgeoisie (\texttt{checks\_bourgeoisie})}

\begin{promptquote}
Does an urban commercial or capital-holding class provide an effective check on the executive
during THIS leader's tenure? An effective check means the body is independent of the executive
and has the de facto capacity to resist or constrain executive action --- not merely formal/de
jure standing. A body that exists on paper but is controlled by the executive does NOT count.
Media, civil society organizations, and ordinary citizens are NOT counted.

Examples: merchants, traders, artisans, financiers, creditors, business/commercial classes.

Coding: 0 = No --- no independent commercial class exists, or it cannot effectively resist the
executive. 1 = Yes --- the commercial class is independent and can, at least on occasion,
effectively resist the executive.
\end{promptquote}

\texttt{checks\_bourgeoisie\_prediction} (int, 0--1):
\begin{itemize}[nosep, leftmargin=1.7em]
  \item 0 = does not effectively check the executive
  \item 1 = effectively checks the executive
\end{itemize}
\texttt{checks\_bourgeoisie\_confidence} (int, 1--100): LLM self-reported confidence score.

\bigskip
\textbf{Checks: Bureaucracy (\texttt{checks\_bureaucracy})}

\begin{promptquote}
Does a professional administrative apparatus provide an effective check on the executive during
THIS leader's tenure? An effective check means the body is independent of the executive and has
the de facto capacity to resist or constrain executive action --- not merely formal/de jure
standing. A body that exists on paper but is controlled by the executive does NOT count. Media,
civil society organizations, and ordinary citizens are NOT counted.

Examples: civil servants, Confucian scholar-officials serving as top-level advisors and
administrators.

Coding: 0 = No --- no independent bureaucracy exists, or it cannot effectively resist the
executive. 1 = Yes --- the bureaucracy is independent and can, at least on occasion, effectively
resist the executive.
\end{promptquote}

\texttt{checks\_bureaucracy\_prediction} (int, 0--1):
\begin{itemize}[nosep, leftmargin=1.7em]
  \item 0 = does not effectively check the executive
  \item 1 = effectively checks the executive
\end{itemize}
\texttt{checks\_bureaucracy\_confidence} (int, 1--100): LLM self-reported confidence score.

\bigskip
\textbf{Checks: Judiciary (\texttt{checks\_judiciary})}

\begin{promptquote}
Do independent adjudicative bodies provide an effective check on the executive during THIS
leader's tenure? An effective check means the body is independent of the executive and has the
de facto capacity to resist or constrain executive action --- not merely formal/de jure standing.
A body that exists on paper but is controlled by the executive does NOT count. Media, civil
society organizations, and ordinary citizens are NOT counted.

Examples: courts of law, tribunals, judicial bodies, legal institutions.

Coding: 0 = No --- no independent judiciary exists, or it cannot effectively resist the
executive. 1 = Yes --- the judiciary is independent and can, at least on occasion, effectively
resist the executive.
\end{promptquote}

\texttt{checks\_judiciary\_prediction} (int, 0--1):
\begin{itemize}[nosep, leftmargin=1.7em]
  \item 0 = does not effectively check the executive
  \item 1 = effectively checks the executive
\end{itemize}
\texttt{checks\_judiciary\_confidence} (int, 1--100): LLM self-reported confidence score.

\bigskip
\textbf{Checks: Assembly (\texttt{checks\_assembly})}

\begin{promptquote}
Does a deliberative representative or popular body provide an effective check on the executive
during THIS leader's tenure? An effective check means the body is independent of the executive
and has the de facto capacity to resist or constrain executive action --- not merely formal/de
jure standing. A body that exists on paper but is controlled by the executive does NOT count.
Media, civil society organizations, and ordinary citizens are NOT counted.

Examples: popular assembly, legislature, parliament.

Coding: 0 = No --- no independent assembly exists, or it cannot effectively resist the executive.
1 = Yes --- an assembly is independent and can, at least on occasion, effectively resist the
executive.
\end{promptquote}

\texttt{checks\_assembly\_prediction} (int, 0--1):
\begin{itemize}[nosep, leftmargin=1.7em]
  \item 0 = does not effectively check the executive
  \item 1 = effectively checks the executive
\end{itemize}
\texttt{checks\_assembly\_confidence} (int, 1--100): LLM self-reported confidence score.

\bigskip
\textbf{Checks: Advisory Council (\texttt{checks\_council})}

\begin{promptquote}
Does an institutionalized elite council adjacent to the executive provide an effective check on
the executive during THIS leader's tenure? An effective check means the body is independent of
the executive and has the de facto capacity to resist or constrain executive action --- not
merely formal/de jure standing. A body that exists on paper but is controlled by the executive
does NOT count. Media, civil society organizations, and ordinary citizens are NOT counted.

Examples: royal council, council of state, regency council, privy council, council of elders.

Coding: 0 = No --- no independent council exists, or it cannot effectively resist the executive.
1 = Yes --- a council is independent and can, at least on occasion, effectively resist the
executive.
\end{promptquote}

\texttt{checks\_council\_prediction} (int, 0--1):
\begin{itemize}[nosep, leftmargin=1.7em]
  \item 0 = does not effectively check the executive
  \item 1 = effectively checks the executive
\end{itemize}
\texttt{checks\_council\_confidence} (int, 1--100): LLM self-reported confidence score.

% ─────────────────────────────────────────────────────────────────────────────
\subsection{Elections (\texttt{elections}) --- Downstream Task}
% ─────────────────────────────────────────────────────────────────────────────

\textbf{Prompt source:} \texttt{pipeline/post\_\allowbreak processing.py}. Run separately, aFTER
the main pipeline. Depends on \texttt{assembly\_\allowbreak prediction}: rows where \texttt{assembly\_\allowbreak prediction}
$\neq$ 2 receive a hard-coded pass-through label of \texttt{"0"} with NO API call;
only rows with \texttt{assembly\_\allowbreak prediction} = 2 get an LLM call.

\begin{promptquote}
You are a political scientist coding executive constraints for historical leaders. Core rule:
focus on THIS specific leader's tenure. When evidence is uncertain: this indicator can safely
default to 0 / None --- if legislative elections existed, the historical record would usually
mention it, so silence indicates absence.

\textbf{Elections.} A legislative election refers to a regularized selection procedure in which
members are chosen by an electorate. That electorate must be considerably larger than the body
itself, though it may be limited to a small portion of the general population. A LARGE ASSEMBLY
that plays some role in policymaking or leadership selection is KNOWN TO EXIST in the polity
below. Your task is to determine whether members of
that assembly were elected to their positions, and --- if so --- whether those elections were
contested by organized factions or parties.

Coding: 0 = No elections --- members of the assembly obtain their seats by appointment,
heredity, or other non-electoral means (e.g., an appointive legislature or polities without a
legislature). 1 = Non-competitive --- elections without identifiable parties or factions, or
where one party/faction monopolizes the field; may consist of a vote by acclamation. 2 =
Competitive --- elections are organized by multiple factions or parties; although the contest
may not be free and fair, contestants have a reasonable expectation of winning power at some
point in the future, and turnover is conceivable.
\end{promptquote}

\medskip
\textit{User prompt template:}
\begin{quote}\small
Classify the elections indicator for the following leader: \\
\textbf{Polity:} \{polity\} \\
\textbf{Leader:} \{name\} \\
\textbf{Tenure:} \{start\_year\}--\{end\_year\}. \\
A LARGE ASSEMBLY that plays some role in policymaking or leadership selection is KNOWN TO EXIST. Return a single
JSON object with all required fields.
\end{quote}

\textbf{Related columns:}
\ra{1.3}
\begin{longtable}{P{4.8cm} P{3cm} P{6.7cm}}
  \toprule
  \thead{Column} & \thead{Type} & \thead{Meaning} \\
  \midrule
  \endhead
  \bottomrule
  \endfoot
  \texttt{elections\_prediction} & int (0/1/2) &
    0 = not elected, or pass-through when \texttt{assembly\_\allowbreak prediction} $\neq$ 2\newline
    1 = elected, no organized factions\newline
    2 = competitive elections (organized factions or parties) \\
  \texttt{elections\_confidence} & int (1--100), null for pass-through rows &
    LLM self-reported confidence score \\
\end{longtable}

% ═══════════════════════════════════════════════════════════════════════════════
\section{Dataset Statistics}
% ═══════════════════════════════════════════════════════════════════════════════

This section reports label distributions and missingness for the current full-dataset snapshot. \texttt{elections} is not yet included
(downstream task, not yet run).

% ─────────────────────────────────────────────────────────────────────────────
\subsection{Missing predictions in this dataset}
% ─────────────────────────────────────────────────────────────────────────────

The \texttt{constitution} task has zero missing predictions (135{,}672 / 135{,}672 parsed). The
\texttt{indicators} task (one combined LLM call per leader-tenure) has a small, uniform number of
missing predictions across almost every indicator --- exactly 14 rows for every indicator except
\texttt{entry} (19) and \texttt{exit} (18). The uniform 14 come from responses whose JSON was
truncated or unparseable for the whole call, which nulls every indicator in that row
simultaneously; \texttt{entry}/\texttt{exit} each have a handful of additional
indicator-specific validation failures.

These rows are already identified with a ready-to-resubmit jsonline file already written. \textbf{As of this
snapshot the retry has not yet been resubmitted}, pending resolution of a temporary Gemini API
access issue. Once resubmitted per the retry-merge protocol described later in this codebook,
these $\leq$19-row gaps should close to zero without touching any other row.

% ─────────────────────────────────────────────────────────────────────────────
\subsection{Label distributions}
% ─────────────────────────────────────────────────────────────────────────────

Counts and percentages of total rows (135{,}672) for every category actually observed, plus the
missing-prediction rows documented above. Category meanings are defined in Section~2.

\ra{1.15}
\begin{longtable}{P{7.6cm} P{2.3cm} P{2.3cm}}
  \toprule
  \thead{Label} & \thead{Count} & \thead{\% of rows} \\
  \midrule
  \endhead
  \bottomrule
  \endfoot
  \multicolumn{3}{l}{\textbf{Constitution (\texttt{constitution})}} \\
  \quad 0 & 59,168 & 43.61\% \\
  \quad 1 & 21,137 & 15.58\% \\
  \quad 2 & 55,367 & 40.81\% \\
  \addlinespace
  \multicolumn{3}{l}{\textbf{Sovereign (\texttt{sovereign})}} \\
  \quad 0 & 58,018 & 42.76\% \\
  \quad 1 & 77,640 & 57.23\% \\
  \quad Missing (unparsed) & 14 & 0.01\% \\
  \addlinespace
  \multicolumn{3}{l}{\textbf{Federalism (\texttt{federalism})}} \\
  \quad 0 & 129,215 & 95.24\% \\
  \quad 1 & 6,443 & 4.75\% \\
  \quad Missing (unparsed) & 14 & 0.01\% \\
  \addlinespace
  \multicolumn{3}{l}{\textbf{Collegiality (\texttt{collegiality})}} \\
  \quad 0 & 117,213 & 86.39\% \\
  \quad 1 & 18,445 & 13.60\% \\
  \quad Missing (unparsed) & 14 & 0.01\% \\
  \addlinespace
  \multicolumn{3}{l}{\textbf{Petition (\texttt{petition})}} \\
  \quad 0 & 81,168 & 59.83\% \\
  \quad 1 & 54,490 & 40.16\% \\
  \quad Missing (unparsed) & 14 & 0.01\% \\
  \addlinespace
  \multicolumn{3}{l}{\textbf{Assembly (\texttt{assembly})}} \\
  \quad 0 & 55,054 & 40.58\% \\
  \quad 1 & 38,555 & 28.42\% \\
  \quad 2 & 38,189 & 28.15\% \\
  \quad 3 & 3,860 & 2.85\% \\
  \quad Missing (unparsed) & 14 & 0.01\% \\
  \addlinespace
  \multicolumn{3}{l}{\textbf{Entry (\texttt{entry})}} \\
  \quad 0 & 8,890 & 6.55\% \\
  \quad 1 & 29,027 & 21.39\% \\
  \quad 2 & 1,719 & 1.27\% \\
  \quad 3 & 6,125 & 4.51\% \\
  \quad 4 & 32,730 & 24.12\% \\
  \quad 5 & 528 & 0.39\% \\
  \quad 6 & 10,664 & 7.86\% \\
  \quad 7 & 19,121 & 14.09\% \\
  \quad 8 & 1,321 & 0.97\% \\
  \quad 9 & 4,154 & 3.06\% \\
  \quad 10 & 7,632 & 5.63\% \\
  \quad 99 & 13,742 & 10.13\% \\
  \quad Missing (unparsed) & 19 & 0.01\% \\
  \addlinespace
  \multicolumn{3}{l}{\textbf{Exit (\texttt{exit})}} \\
  \quad 0 & 4,670 & 3.44\% \\
  \quad 1 & 43,191 & 31.83\% \\
  \quad 2 & 6,308 & 4.65\% \\
  \quad 3 & 137 & 0.10\% \\
  \quad 4 & 532 & 0.39\% \\
  \quad 5 & 34,731 & 25.60\% \\
  \quad 6 & 888 & 0.65\% \\
  \quad 7 & 224 & 0.17\% \\
  \quad 8 & 8,557 & 6.31\% \\
  \quad 9 & 3,578 & 2.64\% \\
  \quad 10 & 790 & 0.58\% \\
  \quad 11 & 2,078 & 1.53\% \\
  \quad 12 & 364 & 0.27\% \\
  \quad 13 & 5,776 & 4.26\% \\
  \quad 14 & 708 & 0.52\% \\
  \quad 99 & 23,122 & 17.04\% \\
  \quad Missing (unparsed) & 18 & 0.01\% \\
  \addlinespace
  \multicolumn{3}{l}{\textbf{Symbolism (\texttt{symbolism})}} \\
  \quad 0 & 35,912 & 26.47\% \\
  \quad 1 & 74,884 & 55.19\% \\
  \quad 2 & 23,401 & 17.25\% \\
  \quad 3 & 1,461 & 1.08\% \\
  \quad Missing (unparsed) & 14 & 0.01\% \\
  \addlinespace
  \multicolumn{3}{l}{\textbf{Checks --- Local / Constituent Units (\texttt{checks\_local})}} \\
  \quad 0 & 109,316 & 80.57\% \\
  \quad 1 & 26,342 & 19.42\% \\
  \quad Missing (unparsed) & 14 & 0.01\% \\
  \addlinespace
  \multicolumn{3}{l}{\textbf{Checks --- Military (\texttt{checks\_military})}} \\
  \quad 0 & 124,971 & 92.11\% \\
  \quad 1 & 10,687 & 7.88\% \\
  \quad Missing (unparsed) & 14 & 0.01\% \\
  \addlinespace
  \multicolumn{3}{l}{\textbf{Checks --- Clergy (\texttt{checks\_clergy})}} \\
  \quad 0 & 113,248 & 83.47\% \\
  \quad 1 & 22,410 & 16.52\% \\
  \quad Missing (unparsed) & 14 & 0.01\% \\
  \addlinespace
  \multicolumn{3}{l}{\textbf{Checks --- Aristocracy (\texttt{checks\_aristocracy})}} \\
  \quad 0 & 97,688 & 72.00\% \\
  \quad 1 & 37,970 & 27.99\% \\
  \quad Missing (unparsed) & 14 & 0.01\% \\
  \addlinespace
  \multicolumn{3}{l}{\textbf{Checks --- Bourgeoisie (\texttt{checks\_bourgeoisie})}} \\
  \quad 0 & 123,999 & 91.40\% \\
  \quad 1 & 11,659 & 8.59\% \\
  \quad Missing (unparsed) & 14 & 0.01\% \\
  \addlinespace
  \multicolumn{3}{l}{\textbf{Checks --- Bureaucracy (\texttt{checks\_bureaucracy})}} \\
  \quad 0 & 129,931 & 95.77\% \\
  \quad 1 & 5,727 & 4.22\% \\
  \quad Missing (unparsed) & 14 & 0.01\% \\
  \addlinespace
  \multicolumn{3}{l}{\textbf{Checks --- Judiciary (\texttt{checks\_judiciary})}} \\
  \quad 0 & 113,783 & 83.87\% \\
  \quad 1 & 21,875 & 16.12\% \\
  \quad Missing (unparsed) & 14 & 0.01\% \\
  \addlinespace
  \multicolumn{3}{l}{\textbf{Checks --- Assembly (\texttt{checks\_assembly})}} \\
  \quad 0 & 99,555 & 73.38\% \\
  \quad 1 & 36,103 & 26.61\% \\
  \quad Missing (unparsed) & 14 & 0.01\% \\
  \addlinespace
  \multicolumn{3}{l}{\textbf{Checks --- Advisory Council (\texttt{checks\_council})}} \\
  \quad 0 & 100,480 & 74.06\% \\
  \quad 1 & 35,178 & 25.93\% \\
  \quad Missing (unparsed) & 14 & 0.01\% \\
\end{longtable}

% ─────────────────────────────────────────────────────────────────────────────
\subsection{Constitution document-metadata missingness}
% ─────────────────────────────────────────────────────────────────────────────

Within the \texttt{constitution} task, the three document-level columns ---
\texttt{constitution\_\allowbreak document\_\allowbreak name}, \texttt{constitution\_\allowbreak
document\_\allowbreak types}, and \texttt{constitution\_\allowbreak year} --- are missing at very
different rates once conditioned on \texttt{constitution\_\allowbreak prediction}:

\ra{1.2}
\begin{longtable}{P{2.4cm} P{2cm} P{2.9cm} P{2.9cm} P{2.9cm}}
  \toprule
  \thead{Prediction} & \thead{Rows} & \thead{document\_name missing} &
  \thead{document\_types missing} & \thead{constitution\_year missing} \\
  \midrule
  \endhead
  \bottomrule
  \endfoot
  0 (no written law) & 59,168 & 59,166 (100.00\%) & 59,166 (100.00\%) & 59,166 (100.00\%) \\
  1 (code of law) & 21,137 & 5 (0.02\%) & 6 (0.03\%) & \textbf{1,519 (7.19\%)} \\
  2 (full constitution) & 55,367 & 0 (0.00\%) & 0 (0.00\%) & \textbf{777 (1.40\%)} \\
\end{longtable}

At prediction = 0 all three columns are missing together, as expected (no document to name, type,
or date). At prediction = 1 or 2, however, \texttt{constitution\_\allowbreak year} is missing far
more often than \texttt{document\_\allowbreak name}/\texttt{document\_\allowbreak types} ---
roughly 300--600$\times$ more rows. This is a deliberate consequence of the prompt wording
(Section~2.1): the model is instructed to give the adoption year \emph{as an exact integer only,
or \texttt{"N/A"} if uncertain --- no approximations (``c.'', ``circa'', ``$\sim$'')}, with no
equivalent exactness requirement on the document's name or type. In practice the model is much
more often able to identify and classify a legal document than to pin down its precise adoption
year, especially for premodern or poorly-documented polities; the asymmetry reflects the model
honestly reporting that uncertainty rather than a parsing defect. (A separate check found 2 rows
with prediction = 0 yet non-null document metadata --- an internally inconsistent LLM output,
negligible at 0.001\% of rows.)

% ═══════════════════════════════════════════════════════════════════════════════
\section{Related Files}
% ═══════════════════════════════════════════════════════════════════════════════

% \subsection{Input files}

% \ra{1.3}
% \begin{longtable}{P{4.2cm} P{10.3cm}}
%   \toprule
%   \thead{File} & \thead{Description} \\
%   \midrule
%   \endhead
%   \bottomrule
%   \endfoot
%   \texttt{data/\allowbreak plt\_\allowbreak dataset.csv} (or \texttt{data/\allowbreak
%   plt\_\allowbreak leaders\_\allowbreak data.csv}) &
%     Leader-level input CSV/JSONL (Section 1). Loaded via
%     \texttt{utils/\allowbreak data\_\allowbreak loader.py::load\_\allowbreak dataframe()}, which
%     accepts both CSV and JSONL. \\
%   \texttt{data/\allowbreak temp/\allowbreak batch\_\{task\}.jsonl} &
%     Batch request file built by \texttt{src/\allowbreak build\_\allowbreak batch\_\allowbreak
%     jsonl.py --task constitution\textbar indicators\textbar elections}. File-upload REST schema
%     (\texttt{\{"key", "request", "metadata"\}}); custom ID format
%     \texttt{"\{row\_idx\}\textbar constitution\textbar\{sc\_idx\}"} or
%     \texttt{"\{row\_idx\}\textbar elections\textbar\{sc\_idx\}"} for those two tasks, and
%     \texttt{"\{row\_idx\}\textbar single\textbar\{sc\_idx\}"} for the combined \texttt{indicators}
%     task (one request covers ALL selected indicators). \texttt{row\_idx} is the 0-based
%     positional index in the input file; \texttt{sc\_idx} = 0 is the initial prediction
%     (\texttt{sc\_idx} $\geq$ 1 only exists when self-consistency sampling is enabled). Split into
%     chunks automatically if a job would exceed the 2\,GB file-upload limit. \\
%   \texttt{data/\allowbreak temp/\allowbreak *\_manifest.json} &
%     Build manifest recording the task, model-agnostic build parameters
%     (\texttt{n\_samples}, \texttt{reasoning}, \texttt{prompt\_version}, \texttt{search\_mode},
%     indicator list), and chunk file list for a given build. \\
% \end{longtable}

\subsection{Output files}

\ra{1.3}
\begin{longtable}{P{4.5cm} P{10cm}}
  \toprule
  \thead{File} & \thead{Description} \\
  \midrule
  \endhead
  \bottomrule
  \endfoot
  \texttt{plt\_constraints.csv} &
    Tabular prediction output. All original input columns are preserved, with the new
    \texttt{\{ind\}\_prediction} / \texttt{\{ind\}\_confidence} (and constitution's three extra)
    columns appended. \\
  \texttt{\{task\}\_\allowbreak provenance.json} &
    Per-run provenance record written by the batch runner: \texttt{model}, \texttt{run\_timestamp},
    \texttt{n\_samples}, \texttt{input\_jsonl} (paths of the request file(s) submitted),
    \texttt{total\_requests}, \texttt{success\_count}, \texttt{response\_failed\_cids} \\
  \texttt{\{task\}\_\allowbreak failed\_\allowbreak
  requests.jsonl} &
    Only written when at least one row failed. Contains exactly the failed rows' requests,
    ready to resubmit as-is to \texttt{jsonl\_\allowbreak batch\_\allowbreak runner.py}. \\
\end{longtable}

% \subsection{Retry-merge protocol}

% To resubmit a \texttt{\_failed\_requests.jsonl}: pass it as \texttt{--input}, point
% \texttt{--dataset} at the PREVIOUS run's \textbf{output} CSV (not the original raw input file),
% and reuse the SAME \texttt{--output} path --- only the retried rows are overwritten and every
% already-successful row is preserved. Merging is strictly positional by row index, so the dataset
% CSV must not be reordered or filtered between rounds. Before submitting, the runner verifies each
% request's embedded \texttt{id} (echoed back from build-time \texttt{row\_data} metadata) against
% the \texttt{--dataset} row at the same position, and aborts on any mismatch, so a
% reordered/filtered dataset fails loudly instead of silently merging predictions onto the wrong
% leader. Repeat until a run produces no \texttt{\_failed\_requests.jsonl}.

\subsection{Chaining the three tasks}

\begin{enumerate}[leftmargin=1.4em]
  \item \texttt{constitution} build + run from the raw input $\rightarrow$ \texttt{plt\_constitution.csv}
  \item \texttt{indicators} build + run from the raw input $\rightarrow$ \texttt{plt\_indicators.csv}
  \item \texttt{elections} build FROM \texttt{plt\_indicators.csv} (requires its
    \texttt{assembly\_prediction} column), run with \texttt{--dataset plt\_indicators.csv} $\rightarrow$
    final merged dataset
\end{enumerate}

\end{document}
